\section*{Capítulo \ref{ch:b3:environmental-toxicology} --- Química ambiental y toxicología}

\begin{solution}{exo:b3:environmental-toxicology:1}
Una molécula diatómica homonuclear tiene una sola vibración, que mantiene nulo el momento dipolar,
y el argón no tiene ninguna vibración: ninguno puede absorber la radiación infrarroja. El $\ce{CO2}$ absorbe gracias a
su modo de flexión (y a su tensión antisimétrica), que crean un dipolo
oscilante; la flexión está en plena emisión de la Tierra.
\end{solution}

\begin{solution}{exo:b3:environmental-toxicology:2}
$10^{0.1} = 1.26$: $[\ce{H+}]$ aumenta un 26\,\%.
\end{solution}

\begin{solution}{exo:b3:environmental-toxicology:3}
$(2.0 + 0.50)\ \unit{mmol/L} \times \qty{100.1}{mg/mmol} = \qty{250}{mg/L}$ de $\ce{CaCO3}$.
\end{solution}

\begin{solution}{exo:b3:environmental-toxicology:4}
N: $0.60/14.0 = \qty{0.043}{mmol/L}$; P: $0.010/31.0 = \qty{3.2e-4}{mmol/L}$; N\,:\,P $= 133$, muy por encima de 16: el fósforo se agota
primero y limita el crecimiento.
\end{solution}

\begin{solution}{exo:b3:environmental-toxicology:5}
$L_0 = \mathrm{DBO_5}/(1 - \eu^{-0.23 \times 5}) = 170/0.683 = \qty{250}{mg/L}$.
\end{solution}

\begin{solution}{exo:b3:environmental-toxicology:6}
$\qty{4.40e-3}{L} \times \qty{0.0200}{mol/L} = \qty{8.80e-5}{mol}$ en \qty{0.1000}{L}: \omterm{def:b3:environmental-toxicology:alkalinity}{alcalinidad} de $\qty{8.80e-4}{mol/L}$; como $\ce{HCO3-}$ (\qty{61.0}{g/mol}),
\qty{53.7}{mg/L}.
\end{solution}

\begin{solution}{exo:b3:environmental-toxicology:7}
$K_d = 0.020 \times 300 = \qty{6.0}{L/kg}$. Sorbido/total $= K_dm_s/(K_dm_s + V_w) = 9.0/(9.0 + 0.30) = 0.97$: un 97\,\% sorbido.
\end{solution}

\begin{solution}{exo:b3:environmental-toxicology:8}
Benceno: $\log\mathrm{FBC} = 0.85 \times 2.13 - 0.70 = 1.11$, FBC de unos 13. PCB: $0.85 \times 6.67 - 0.70 = 4.97$, FBC de unos
\num{9e4}: se concentra en los peces unas siete mil veces más que el benceno; con un $K_{ow}$ así,
la regresión está cerca de su límite de validez.
\end{solution}

\begin{solution}{exo:b3:environmental-toxicology:9}
$\qty{192}{mg/kg} \times \qty{70}{kg} = \qty{13}{g}$, unas 150 tazas. Las especies difieren en absorción y en metabolismo,
y una $\mathrm{DL_{50}}$ no es un umbral seguro; los efectos nocivos empiezan muy por debajo de ella.
\end{solution}

\begin{solution}{exo:b3:environmental-toxicology:10}
PNEC $= \qty{0.50}{mg/L}/1000 = \qty{0.50}{\micro g/L}$; $\mathrm{CR} = 0.20/0.50 = 0.40$, inferior a 1: sin motivo de preocupación con esta exposición.
\end{solution}

\begin{solution}{exo:b3:environmental-toxicology:11}
Plancton/agua: $\qty{0.015}{mg/kg}/\qty{2.0e-6}{mg/L} = \qty{7.5e3}{L/kg}$ (bioconcentración). Peces/plancton: $0.30/0.015 = 20$;
ave/peces: $6.0/0.30 = 20$ (\omterm{def:b3:environmental-toxicology:bio}{biomagnificación} en cada eslabón); ave/agua: $\num{3e6}$.
\end{solution}

\begin{solution}{exo:b3:environmental-toxicology:12}
$\qty{1.0e9}{g}/\qty{64.1}{g/mol} = \qty{1.56e7}{mol}$ de $\ce{SO2}$, que dan otros tantos moles de $\ce{H2SO4}$: $\qty{1.53e9}{g}$, unas \qty{1500}{t}. Este
libera $\qty{3.12e7}{mol}$ de $\ce{H+}$; a pH 4.0, $[\ce{H+}] = \qty{1.0e-4}{mol/L}$: $\qty{3.1e11}{L}$ de lluvia.
\end{solution}

\begin{solution}{pb:b3:environmental-toxicology:1}
\textbf{1.} Es diez mil veces más soluble en octanol que en agua: hidrófobo,
se sorberá en la materia orgánica y se acumulará en la grasa.

\textbf{2.} $K_{oc} = 0.41 \times 10^4 = \qty{4100}{L/kg}$; $K_{sw} = f_{oc}K_{oc} = 0.040 \times 4100 = \qty{164}{L/kg}$.

\textbf{3.} $K_{fw} = 0.048 \times 10^4 = \qty{480}{L/kg}$.

\textbf{4.} No: en el equilibrio, el aire contiene $10^{-4}$ veces la concentración del agua.

\textbf{5.} La molécula neutra e hidrófoba se disuelve en la materia orgánica, como en el
octanol; las superficies minerales son polares y están cubiertas de agua.

\textbf{6.} El sedimento, cuya capacidad, $K_{sw}m_s$, es la mayor.

\textbf{7.} $V_w = \qty{1.0e9}{L}$; $K_{aw}V_a = \qty{1.0e8}{L}$; $K_{sw}m_s = \qty{1.64e10}{L}$; $K_{fw}m_f = \qty{4.8e6}{L}$.

\textbf{8.} $C_w = \qty{1.0e8}{mg}/\qty{1.75e10}{L} = \qty{5.7e-3}{mg/L}$, \qty{5.7}{\micro g/L}.

\textbf{9.} Agua, \qty{5.7}{kg}; aire, \qty{0.57}{kg}; sedimento, \qty{94}{kg}; peces, \qty{0.027}{kg}.

\textbf{10.} Aire, $\qty{5.7e-7}{mg/L}$; sedimento, \qty{0.94}{mg/kg}; peces, \qty{2.7}{mg/kg}.

\textbf{11.} El equilibrio entre todos los compartimentos a la vez, sin degradación, sin
entradas ni salidas, y con compartimentos bien mezclados.

\textbf{12.} Con degradación y flujos en estado estacionario (nivel II), el inventario lo fijan
las velocidades de entrada y de pérdida; con velocidades de transferencia entre compartimentos (nivel III),
las concentraciones ya no están en razones de equilibrio, y el compartimento que
recibe la entrada se enriquece.

\textbf{13.} $k = \ln 2/\qty{60}{d} = \qty{0.0116}{d^{-1}}$.

\textbf{14.} $\eu^{-0.0116 \times 365} = 0.0145$: queda un 1.5\,\%.

\textbf{15.} $\qty{5.7}{\micro g/L} \times 0.0145 = \qty{0.083}{\micro g/L}$.

\textbf{16.} La degradación suele ser más lenta en un sedimento frío y anóxico, y el plaguicida
sorbido vuelve lentamente al agua, lo que lo mantiene presente después de que el
agua se haya limpiado.

\textbf{17.} Enlaces que resisten la hidrólisis, la oxidación y el ataque microbiano (C--Cl
aromático, por ejemplo), una baja solubilidad en agua y una sorción que lo protege de los
degradadores.

\textbf{18.} PNEC $= \qty{50}{\micro g/L}/10 = \qty{5.0}{\micro g/L}$.

\textbf{19.} $\mathrm{CR} = 5.7/5.0 = 1.1$.

\textbf{20.} Justo por encima de 1: se indica un riesgo, y hay que afinar la evaluación (mejores
datos de toxicidad, concentraciones medidas) o reducir la exposición.

\textbf{21.} \qty{2.7}{mg/kg}, casi el triple del umbral de \qty{1.0}{mg/kg}: los peces no deben consumirse
hasta que la concentración disminuya.

\textbf{22.} Las concentraciones en el agua, el sedimento y los peces a lo largo del tiempo y en varios puntos,
para comprobar el modelo y seguir la disminución.

\textbf{23.} Menos plaguicida aplicado, aplicación lejos de la lluvia y de la orilla,
franjas de protección que retengan la escorrentía, o una alternativa menos persistente.

\textbf{24.} Divide el nivel sin efecto medido para cubrir la distancia entre unas pocas
especies de laboratorio y todo un ecosistema; un factor de 10 en lugar de 1000 cambia la
PNEC, y el veredicto, en un factor cien.

\textbf{25.} El \omterm{def:b3:environmental-toxicology:ecotoxicology}{cociente de riesgo} PEC/PNEC es de alrededor de 1.1, justo por encima de 1.
\end{solution}
